%***************************************************************************************************
% preamble/pagestyles.tex - nagłówki i stopki (fancyhdr).
% Style nazwane (zamiast nadpisywania stylu "plain" w trakcie dokumentu):
%   frontstyle    - czysta strona (strona tytułowa, spis treści, streszczenie, spis skrótów)
%   chapterfirst  - pierwsza strona rozdziału: tylko stopka (numer strony i przedmiot)
%   mainstyle     - rozdziały i załączniki: nagłówek "Rozdział N" / "Załącznik A" + stopka
%   bibstyle, lofstyle, lotstyle, lolstyle, loastyle - bibliografia i spisy (nagłówek z tytułem spisu)
% Nagłówek nie używa współrzędnych \put - dopasowuje się do szerokości tekstu.
%***************************************************************************************************

% Etykieta rozdziału: "Rozdział 3" w rozdziałach, "Załącznik A" po \appendix
\makeatletter
\newcommand{\HeaderChapterLabel}{\@chapapp\ \thechapter}
\makeatother

% Wiersz z nazwą uczelni i wydziału nad linią nagłówka
\newcommand{\HeaderInstitution}{\footnotesize\color{dgray}\University\ \textendash\ \Faculty}

% Strony nieparzyste: nazwa uczelni do lewej, linia i etykieta do prawej
\newcommand{\HeaderOddBox}[1]{%
  \parbox[b]{\headwidth}{%
    \raggedright\noindent\HeaderInstitution\par\nointerlineskip
    \noindent\makebox[\headwidth]{{\color{Accent}\hrulefill}\hspace{1em}\large #1}}}

% Strony parzyste: lustrzane odbicie
\newcommand{\HeaderEvenBox}[1]{%
  \parbox[b]{\headwidth}{%
    \raggedleft\noindent\HeaderInstitution\par\nointerlineskip
    \noindent\makebox[\headwidth]{\large #1\hspace{1em}{\color{Accent}\hrulefill}}}}

% Stopka: numer strony przy zewnętrznej krawędzi, przedmiot przy wewnętrznej, cienka linia nad nią
\newcommand{\FooterSetup}{%
  \fancyhf{}%
  \fancyfoot[LE,RO]{\thepage}%
  \fancyfoot[LO,RE]{\small\Subject}%
  \renewcommand{\headrulewidth}{0pt}%
  \renewcommand{\footrulewidth}{0.2pt}%
}

% Czysta strona (bez nagłówka, numeru i stopki)
\fancypagestyle{frontstyle}{%
  \fancyhf{}%
  \renewcommand{\headrulewidth}{0pt}%
  \renewcommand{\footrulewidth}{0pt}%
}

% Pierwsza strona rozdziału (przypisywana do \chapter w main.tex);
% "plain" zdefiniowany raz, tak samo, na wypadek użycia przez inne polecenia
\fancypagestyle{chapterfirst}{\FooterSetup}
\fancypagestyle{plain}{\FooterSetup}

% Generator stylu z nagłówkiem: \DefineHeaderStyle{nazwa}{tekst etykiety w nagłówku}
\newcommand{\DefineHeaderStyle}[2]{%
  \fancypagestyle{#1}{%
    \FooterSetup
    \fancyhead[RO]{\HeaderOddBox{#2}}%
    \fancyhead[LE]{\HeaderEvenBox{#2}}%
  }%
}

\DefineHeaderStyle{mainstyle}{\HeaderChapterLabel}
\DefineHeaderStyle{bibstyle}{\bibname}
\DefineHeaderStyle{lofstyle}{\listfigurename}
\DefineHeaderStyle{lotstyle}{\listtablename}
\DefineHeaderStyle{lolstyle}{\lstlistlistingname}
\DefineHeaderStyle{loastyle}{\listalgorithmcfname}

% Wstawia spis (nagłówek, wpis w spisie treści, właściwy styl stron):
%   \BackList{styl}{tytuł wpisu w spisie treści}{polecenie spisu}
\newcommand{\BackList}[3]{%
  \clearpage\phantomsection
  \pagestyle{#1}%
  \addcontentsline{toc}{chapter}{#2}%
  #3%
}
